\chapter{供应链、规模化与单位经济性}
\label{chap:supply-chain-unit-economics}

\section{BOM 只是总成本的起点}

机器人交付同时消耗本体、末端、工装、集成、验证、运维、人工恢复和停机机会成本。供应商未披露的 BOM 不应被猜测；更可复算的对象是特定任务在合同期内的总拥有成本（TCO）：
\begin{equation}
C_{\mathrm{TCO}}=C_{\mathrm{cap}}+C_{\mathrm{integration}}+C_{\mathrm{maintenance}}+C_{\mathrm{energy}}+C_{\mathrm{human}}+C_{\mathrm{downtime}}-V_{\mathrm{residual}}.
\label{eq:robot-tco}
\end{equation}
若合同期合格产出为 $N_{\mathrm{good}}$，任务单位成本为 $c_u=C_{\mathrm{TCO}}/N_{\mathrm{good}}$。其中资本项需按寿命和资金成本摊销；人工项必须包含监督、异常清除、远程接管和维护，而不只计算上下料员。

制造 KPI 应有统一定义。ISO 22400-1 给出制造运营 KPI 的通用框架\cite{iso22400_1_2014}。本书采用
\begin{equation}
N_{\mathrm{good}}=H_s\,A\,r\,Y,
\qquad A=\frac{H_s-H_{\mathrm{down}}}{H_s},
\label{eq:effective-output}
\end{equation}
其中 $H_s$ 为排程小时，$A$ 为可用率，$r$ 为运行时理论产率，$Y$ 为良率。提高峰值动作速度只改变 $r$；若跌倒、热保护、换电和等待恢复拉低 $A$，总产出仍可能下降。

\section{关键部件的耦合约束}

执行器把电机、传动、轴承、编码器、驱动器、制动和热路径耦合在关节内。峰值扭矩不能代表连续能力；减速器间隙、效率、冲击寿命与反驱性共同决定控制品质。灵巧手把执行器密度、线缆寿命、触觉、污染防护和维修性压缩到更小体积，通常也是长尾故障源。

感知与算力同样受功耗和散热约束。以 Jetson AGX Orin 为例，官方规格给出 15--60 W 可配置功耗区间\cite{nvidia2026orin}；这只是计算模块而非整机功耗。提高视觉分辨率、模型频率或冗余感知会增加计算、内存带宽、热设计和电池负担，并可能减少有效工作时间。

\begin{table}[H]
\begingroup\hbadness=10000
\centering
\caption{供应链对象必须连到系统指标}
\label{tab:supply-chain-couplings}
\begin{tabularx}{\textwidth}{p{0.16\textwidth}YYY}
\toprule
对象 & 采购规格不够的原因 & 系统级验证 & 常见规模化失效 \\
\midrule
执行器/减速器 & 峰值不等于热稳态与寿命 & 工况谱、冲击、回差、热循环 & 批次离散、密封、润滑、线缆 \\
传感器 & 精度不等于污染/遮挡鲁棒性 & 标定漂移、同步、失效注入 & 供应替代导致重新标定 \\
算力/电池 & TOPS、kWh 不等于任务续航 & 端到端时延、峰值电流、热降频 & 模块 EOL、认证和热一致性 \\
灵巧手/末端 & 自由度不等于可维护操作能力 & 抓取谱、碰撞、清洁与换件 & 小件多、维修慢、寿命分散 \\
\bottomrule
\end{tabularx}
\endgroup
\end{table}

Figure 的 BotQ 材料说明其从 CNC 原型件转向压铸、注塑、冲压等模具工艺，并建立 MES、在制检测和供应商准入\cite{figure2025botq,figure2026production}。这说明“量产设计”改变零件数、连接方式、测试点和工艺，而不是把原型图纸简单复制；但产能、产量、首过良率和客户良率仍是四个不同指标。

\section{一个可复算的任务级算例}

设某三年任务单元的资本与集成摊销为 90 万元，运维和能源 30 万元，排程 12,000 小时，理论产率 50 件/小时，良率 $Y=0.98$。表\ref{tab:unit-economics-sensitivity}仅为教学假设，不代表任何企业报价。人工恢复按 80 元/小时、停机机会损失按 600 元/小时计算。

\begin{table}[H]
\centering
\caption{可用率对单位成本的敏感性（教学算例）}
\label{tab:unit-economics-sensitivity}
\begin{tabular}{lrrr}
\toprule
情景 & 可用率 $A$ & 合格产出/件 & 基础摊销/元每件 \\
\midrule
调试态 & 0.60 & 352,800 & 3.40 \\
稳定态 & 0.85 & 499,800 & 2.40 \\
高可用 & 0.95 & 558,600 & 2.15 \\
\bottomrule
\end{tabular}
\end{table}

基础摊销仅用 120 万元除以产出，尚未计人工和停机。若调试态比稳定态多停机 3,000 小时，停机损失增加 180 万元，远大于假设中的能源项。因此采购评审需要对 $A$、$Y$、人工介入率和故障修复时间做敏感性分析；用最低硬件售价做决定会把主要成本留在合同外。

\section{销售、集成、RaaS 与软件平台}

硬件销售把资产和利用率风险更多交给客户，适合任务稳定且客户有运维能力的场景；系统集成适合一次性工位改造，但收入与工程人力耦合；RaaS 把前期资本转换为服务费，供应商则承担 uptime、备件和残值风险。Agility--GXO 的多年期安排证明 RaaS 已进入人形机器人合同实践\cite{agility2024gxo}，却不自动证明其服务毛利为正。

软件/数据平台只有在多台、多站点复用时才可能脱离项目制。商业合同至少应定义任务和验收、可用率窗口、计划停机、人工责任、数据权利、软件更新、网络安全、事故责任、退出和残值；否则“按结果付费”无法计算结果。

\section{知识小结与综述判断}

供应链的技术问题是部件性能、工艺能力、批次一致性、可测试性和可维修性的联合优化；商业问题是把这些变量折算为合格产出成本。综述判断是：规模化的关键不在某个 BOM 数字，而在设计降复杂度、制造良率、现场可用率和低人工介入共同收敛。融资或规划产能只能说明资源，不能替代单位经济性。

